% authority: science-draft
% status: draft/control
% route: ontology-law-research-packet
% audit_result: source_pure_textually_formally_insufficient_repair_required
% adoption_status: blocked_adoption_open_continuation
\documentclass[11pt]{article}
\usepackage[margin=1in]{geometry}
\usepackage{amsmath,amssymb,amsthm}
\usepackage{booktabs}
\usepackage{longtable}
\usepackage{enumitem}

\newtheorem{lemma}{Lemma}
\newtheorem{proposition}{Proposition}
\newcommand{\Cert}{\mathsf{Cert}_{\mathrm{src}}}
\newcommand{\Accept}{\mathsf{Accept}_{\mathrm{src}}}
\newcommand{\Law}{\mathsf{EqSrcClosureLaw}^{\mathrm{cand}}_{\mathrm{src}}}
\newcommand{\Block}{\mathsf{Block}}

\title{EqSrc Closure Source-Law Candidate Smuggling Audit v1}
\author{Smuggling Auditor control packet RT-20260718-017}
\date{2026-07-19}

\begin{document}
\maketitle

\section{Control Status}

This draft/control artifact audits the proposal-only candidate
\(\Law\) in
\texttt{research\_control/tasks/RT-20260718-015/artifacts/
eqsrc\_family\_closure\_source\_law\_candidate\_v1.tex}.
The candidate remains source-extension data with adoption status
\texttt{blocked\_adoption\_open\_continuation}.

\begin{verbatim}
audit_result: source_pure_textually_formally_insufficient_repair_required
source_purity_result: pass_with_semantic_guard
formal_sufficiency_result: fail_totality_not_stated
obstruction_id: OBST-EQSRC-CLOSURE-TOTALITY-001
next_route: bounded_ontology_formalizer_candidate_repair
general_EqSrc_discharged: false
RetainH_adopted: false
GenH_adopted: false
distance_to_gr_delta: obstruction_recorded_no_ledger_change
\end{verbatim}

The audit finds no explicit target atlas, target metric, detector, benchmark,
registry, role, validator, generated derivative, handoff, or checkpoint
premise in the candidate. That textual source-purity result does not rescue the
conditional proposition: the written clauses do not assert the accepted
identity, inverse, and composition existence/totality used by its proof.

\section{Artifact and Question Under Audit}

The relation under audit is
\[
 A\sim_{\mathrm{cand}}B
 \quad\Longleftrightarrow\quad
 \exists w\in\Cert(A,B)\ \Accept(w).
\]
The candidate then claims that satisfying every candidate-law clause makes
\(\sim_{\mathrm{cand}}\) an equivalence relation.

The exact audit question is twofold:
\begin{enumerate}[leftmargin=*]
  \item Are the vocabulary and proof premises source-defined independently of
  target or process success?
  \item Do the clauses, read as written, entail the accepted witnesses needed
  for reflexivity, symmetry, and transitivity?
\end{enumerate}
These questions are independent. A source-pure statement can still be
formally insufficient.

\section{Audit Method}

The audit separates four statuses:
\begin{itemize}[leftmargin=*]
  \item \texttt{not\_detected}: no written premise uses the inspected
  smuggling channel;
  \item \texttt{semantic\_guard\_required}: a source-labelled term is not yet
  closed by an explicit source-only grammar;
  \item \texttt{formal\_repair\_required}: the proof uses a stronger
  quantifier or operation property than the clauses state;
  \item \texttt{fails\_closed}: a forbidden target or process premise is
  actually required.
\end{itemize}

``May be accepted only if \(P\)'' is read as the necessary condition
\(\Accept(x)\Rightarrow P(x)\). It is not read as either existence of \(x\) or
the sufficient condition \(P(x)\Rightarrow\Accept(x)\). This standard
quantifier discipline prevents a certificate policy from being laundered into
a closure theorem.

\section{Source-Purity Audit}

The candidate explicitly excludes target atlas, target metric, detector
outcome, benchmark result, generated derivative, registry row, role
declaration, and validator result from certificate evidence. Its no-target
grammar also requires definitions and witnesses to precede target
reconstruction. No written proof step invokes any excluded object.

\begin{longtable}{p{0.31\linewidth}p{0.18\linewidth}p{0.43\linewidth}}
\toprule
Audit channel & Finding & Reason \\
\midrule
target topology or atlas & not detected & No chart, transition map, target topology, or target differentiability predicate is used. \\
target Lorentzian metric & not detected & No metric, connection, curvature, cone, geodesic, or proper-time normalization supplies a certificate. \\
detector or empirical protocol & not detected & No detector outcome, calibration table, or empirical acceptance rule is a premise. \\
benchmark behavior & not detected & Exact-GR recovery or benchmark success does not create a source certificate. \\
matter or field equations & not detected & No stress-energy, matter action, Einstein equation, or downstream metric result is used. \\
registry or generated derivative & not detected & Registry and wiki state identify artifacts only; they do not create mathematical witnesses. \\
role, handoff, or checkpoint & not detected & Execution authority and provenance do not entail certificate existence or acceptance. \\
validator result & not detected & A validation receipt checks artifact structure; it cannot make a source law true. \\
\bottomrule
\end{longtable}

\paragraph{Semantic guard.}
The phrases ``source-invariant ledger annotations,'' ``permitted proxy incidence,''
``negative controls,'' and operation ``provenance'' are labelled
source-side but not given closed domains and decision rules in the candidate.
No target import is detected as written. A repair must nevertheless define
these terms entirely in the source grammar, or explicitly bind them to
already-registered source primitives, so that later implementations cannot
populate them using target success or process status.

\section{Formal Sufficiency Audit}

\subsection{Identity}

The identity clause asserts
\(\mathbf 1_A\in\Cert(A,A)\), but does not state
\(\Accept(\mathbf 1_A)\). Fail-closed acceptance gives necessary conditions
for acceptance and does not assert that the identity satisfies them.
Therefore the proof phrase ``accepted identity token'' is stronger than the
written clause.

\subsection{Inverse}

For \(w\in\Cert(A,B)\), the inverse clause says that
\(\operatorname{Inv}(w;\pi_{\mathrm{inv}})\) may be accepted only if its
provenance certifies endpoint reversal and it lies in \(\Cert(B,A)\).
This constrains any accepted inverse token. It does not assert that, for every
accepted \(w\), an inverse provenance token exists, passes the guards, and is
accepted. Therefore symmetry does not follow.

\subsection{Composition}

Likewise the composition clause constrains when a composite token may be
accepted. It does not assert that every composable pair of accepted
certificates has an accepted composite in \(\Cert(A,C)\). Therefore
transitivity does not follow.

\section{Minimal Source-Side Countermodels}

\begin{proposition}[Reflexivity countermodel]
There is a source-only interpretation satisfying the written identity and
acceptance-policy clauses for which \(\sim_{\mathrm{cand}}\) is not reflexive.
\end{proposition}

\begin{proof}
Let the family be \(\{A\}\), let
\(\Cert(A,A)=\{\mathbf 1_A\}\), and set
\(\Accept(\mathbf 1_A)=\mathrm{false}\). The typed identity exists. Every
fail-closed statement of the form ``accepted only when'' is satisfied
vacuously. Yet there is no accepted certificate \(A\to A\), so
\(A\not\sim_{\mathrm{cand}}A\).
\end{proof}

\begin{proposition}[Symmetry countermodel]
There is a source-only interpretation satisfying the written one-way inverse
admissibility condition for which \(\sim_{\mathrm{cand}}\) is not symmetric.
\end{proposition}

\begin{proof}
Let the family be \(\{A,B\}\). Give both objects accepted identity tokens,
take \(\Cert(A,B)=\{w\}\) with \(\Accept(w)=\mathrm{true}\), and take
\(\Cert(B,A)=\varnothing\). Supply neutral, source-only ledger annotations.
No inverse token is accepted, so the implication constraining accepted inverse
tokens is true vacuously. Then \(A\sim_{\mathrm{cand}}B\) but
\(B\not\sim_{\mathrm{cand}}A\).
\end{proof}

\begin{proposition}[Transitivity countermodel]
There is a source-only interpretation satisfying the written one-way
composition admissibility condition for which \(\sim_{\mathrm{cand}}\) is not
transitive.
\end{proposition}

\begin{proof}
Let the family be \(\{A,B,C\}\), give all objects accepted identities, and give
accepted certificates \(w:A\to B\) and \(v:B\to C\), but no accepted
certificate \(A\to C\). Accept no composite token. The restriction on any
accepted composite is true vacuously. Thus \(A\sim_{\mathrm{cand}}B\) and
\(B\sim_{\mathrm{cand}}C\), while
\(A\not\sim_{\mathrm{cand}}C\).
\end{proof}

These finite models use no target atlas, target metric, detector, benchmark, or
process authority. They identify a local statement-level obstruction, not a
global no-go theorem.

\section{Exact Obstruction}

\begin{verbatim}
obstruction_id: OBST-EQSRC-CLOSURE-TOTALITY-001
scope: source_extension_candidate
failed_object: conditional_equivalence_proposition_as_written
exact_failure: admissibility_filters_do_not_entail_accepted_closure_witnesses
current_ontology_implication: does_not_derive
source_extension_implication: repair_allowed
consequence: repair_candidate_allowed
\end{verbatim}

The candidate proof imports no target datum, but it does import mathematical
strength not stated in the definition. Its ``accepted'' identity, inverse, and
composition operations are treated as total on the relevant accepted
certificates. The definition supplies at most existence of an unaccepted
identity token plus necessary conditions on any inverse or composite that
happens to be accepted.

\section{Required Repair}

Before Refuter stress, a bounded Ontology Formalizer repair should state all
of the following explicitly:
\begin{enumerate}[leftmargin=*]
  \item for every \(A\), \(\Accept(\mathbf 1_A)\);
  \item for every accepted \(w:A\to B\), existence of source-only
  \(\pi_{\mathrm{inv}}\) for which
  \(\Accept(\operatorname{Inv}(w;\pi_{\mathrm{inv}}))\);
  \item for every composable accepted \(w:A\to B\) and \(v:B\to C\),
  existence of source-only \(\pi_{\mathrm{comp}}\) for which
  \(\Accept(\operatorname{Comp}(v,w;\pi_{\mathrm{comp}}))\);
  \item source-only totality domains for ledger congruence, negative controls,
  proxy incidence, and provenance;
  \item the existing endpoint, involution, associativity, fail-closed, and
  no-target guards as separate soundness conditions.
\end{enumerate}

These are substantive source-extension axioms, not consequences of a
validator or of current ontology. The repair may remain proposal-only; any
adoption stays human-gated.

\section{Audit Result}

The candidate is \texttt{source\_pure\_as\_written} with a
\texttt{semantic\_guard\_required} for its opaque source-labelled terms.
It is simultaneously
\texttt{formally\_insufficient\_as\_written}. The exact result is
\[
 \boxed{\texttt{repair\_required\_before\_refuter\_stress}}.
\]

This result supersedes neither the earlier conditional H1--H7 audit nor its
source-purity finding. The earlier theorem explicitly supplied closure
hypotheses. The present candidate attempts to package those hypotheses as a
law interface, but its modal wording does not yet state their totality.

\section{Distance-to-GR Effect}

\begin{longtable}{p{0.32\linewidth}p{0.54\linewidth}}
\toprule
Burden & Status after audit \\
\midrule
Source ontology primitives & Unchanged; no canonical ontology edit or adoption. \\
Source equivalence EqSrc & A precise proposal-level totality obstruction is recorded; general EqSrc remains undischarged. \\
RetainH and GenH & Unchanged and not adopted. \\
\(M_{\mathrm{src}}\) and \(g_{\mathrm{eff}}\) & Unchanged; no source-manifold or metric result. \\
Matter coupling and Einstein equations & Blocked and unchanged. \\
Finite-variation robustness & Unchanged; the local countermodels are not a robustness theorem. \\
Benchmark and Gate Chair & Blocked; no promotion or verdict. \\
Route status & Not frozen; one exact proposal-only repair is available before stress. \\
\bottomrule
\end{longtable}

The audit sharpens the obstruction record but does not update the
Distance-to-GR ledger and unlocks no downstream physics object.

\section{Forbidden Conclusions}

This audit does not establish:
\begin{itemize}[leftmargin=*]
  \item canonical ontology modification or source-law adoption;
  \item general EqSrc discharge;
  \item RetainH or GenH adoption;
  \item target atlas, metric, detector, or benchmark authority;
  \item \(M_{\mathrm{src}}\), \(g_{\mathrm{eff}}\), matter coupling, or
  Einstein equations;
  \item benchmark promotion or a Gate Chair verdict;
  \item completed derivation;
  \item impossibility of a conservative repair, future source-extension
  impossibility, or global theory rejection.
\end{itemize}

\section{Next Route}

The next lawful route is one bounded
\texttt{ontology-formalizer@0.2.0} proposal-only repair of the same candidate.
It must add accepted-existence/totality clauses and close the opaque
source-grammar terms. Refuter stress follows only after that repaired artifact
passes a fresh Smuggling Auditor check. No human gate is required merely to
prepare or audit the repair; adoption remains human-gated.

\section{Source Materials}

The AEther-Flow Research Project. (2026a). \emph{Proposal-only source-law
candidate for EqSrc family closure} [Research-control TeX artifact].
\texttt{research\_control/tasks/RT-20260718-015/artifacts/
eqsrc\_family\_closure\_source\_law\_candidate\_v1.tex}.

The AEther-Flow Research Project. (2026b). \emph{EqSrc family-closure
smuggling audit v1} [Research-control TeX artifact].
\texttt{research\_control/tasks/RT-20260707-022/artifacts/
eqsrc\_family\_closure\_smuggling\_audit\_v1.tex}.

The AEther-Flow Research Project. (2026c). \emph{Handoff 0741}
[Research-control handoff].
\texttt{research\_control/handoffs/handoff-0741.yaml}.

\end{document}
